\section{Course terms and textbook numbers}\label{app:terms}
% claims used here: C1 C2 C12-C19 C21 C22 C25 C26 C30 C31 m2-1 m2-2 m3-1 m3-2

This appendix lists every term of the course that the paper uses, with the number and
printed page in the course textbook~\cite{petrosyan2015grafneri} and the English word. Page numbers are the
printed page numbers of the textbook; the PDF page is one higher. The textbook's own word is given in
Latin letters as ``path'', ``simple path'' and so on; its Armenian words are not reproduced.
Terms that the pages checked give no number for are marked ``none''.

\begin{table}[ht]
\centering
\caption{Course terms, textbook numbers and pages, English words, and where they are defined in this paper.}
\label{tab:terms}
\small
\begin{tabular}{@{}p{3.2cm}p{3.0cm}p{1.2cm}p{3.0cm}p{3.0cm}@{}}
\toprule
Term in this paper & Textbook number & Printed page & Textbook word if different & Defined in \\
\midrule
simple undirected graph & Definition 1.1.1 & 9 & graph & \cref{def:graph} \\
adjacent & Definition 1.1.4 & 11 & --- & \cref{def:adjacent} \\
adjacency matrix & none (after Def.\ 1.1.6) & 11 & $A(G)$ & \cref{def:adjmatrix} \\
degree & Definition 1.2.1 & 16 & $d_G(v)$, $d(v)$ & \cref{def:degree} \\
sum of degrees & Theorem 1.2.1 (Euler) & 16 & --- & \cref{prop:rowsums} \\
walk & Definition 1.2.10 & 21 & walk & \cref{def:walk} \\
closed walk & Definition 1.2.11 & 22 & closed walk & \cref{def:closedwalk} \\
trail & Definition 1.2.12 & 22 & ``path'' & \cref{def:trail} \\
path & Definition 1.2.13 & 22 & ``simple path'' & \cref{def:path} \\
cycle & Definition 1.2.14 & 22 & ``cycle'' or ``closed path'' & \cref{def:cycle} \\
simple cycle & Definition 1.2.15 & 22 & ``simple cycle'' & \cref{def:simplecycle} \\
distance & Definition 1.2.16 & 22--23 & distance & \cref{def:distance} \\
walk contains a path & Lemma 1.2.1 (for $u\ne v$) & 23 & --- & \cref{lem:walkpath} \\
connected & Definition 2.1.1 & 48 & connected & \cref{def:connected} \\
connected component & none (prose) & 49 & connectivity component & \cref{def:component} \\
connected iff one component & Remark 2.1.1 & 49 & --- & \cref{rem:component} \\
path between $n_i$ and $n_j$ & Definition 1.2.13 & 22 & ``simple path'' & \cref{def:pathbetween} \\
\bottomrule
\end{tabular}
\end{table}

\paragraph{Figures of the textbook used.}
Fig.~1.2.5 (printed p.~22) is the source of the examples for walks, trails, paths and cycles; \cref{fig:textbook} draws the seven edges that those examples force. Fig.~2.1.1 (p.~48) shows three connected graphs and Fig.~2.1.2 (p.~49) a graph with four components~\cite{petrosyan2015grafneri}.

\paragraph{A note on print order.}
In the printed text, the textbook's example sentence for Definition~1.2.13 (the sequence $v_1,v_3,v_2,v_4,v_3,v_5$ is not a simple path, while $v_1,v_3,v_4,v_5$ is) stands after Definition~1.2.14 and not directly after Definition~1.2.13; its content is about Definition~1.2.13, and this paper cites it by page (p.~22), not by position.

\paragraph{Results that the textbook does not contain.}
The walk-counting theorem (\cref{thm:walkcount}) and the matrix test for connectivity (\cref{thm:test}) are not in the pages checked and were not found by a full-text search of the book (\cref{sec:provenance}), so their absence from the book is \kindtag{qualified}. For the walk-counting theorem the paper cites the University of Galway notes, and for the statement of the test Paramadevan and Sotheeswaran and Weisstein; the proofs are the paper's own.

\paragraph{The board's letters.}
The board writes $(G,V)$ with $V$ the edge set, $e_{ij}\in V$, $|G|=N$ and $n_i$~\cite{whiteboard1,whiteboard2}; the textbook writes $G=(V,E)$, $e\in E$, $n$ and $v_i$. The paper uses $G=(V,E)$, $N$ and $n_i$ (\cref{tab:notation}).
